\documentclass[12pt,a4paper]{article}

%% ------------------------------------------------------------------
%% Packages
%% ------------------------------------------------------------------
\usepackage[utf8]{inputenc}
\usepackage[T1]{fontenc}
\usepackage[english]{babel}
\usepackage{geometry}
\usepackage{amsmath}
\usepackage{amssymb}
\usepackage{booktabs}
\usepackage{array}
\usepackage{longtable}
\usepackage{graphicx}
\usepackage{float}
\usepackage{caption}
\usepackage{subcaption}
\usepackage{hyperref}
\usepackage{xcolor}
\usepackage{listings}
\usepackage{parskip}
\usepackage{enumitem}
\usepackage{titlesec}
\usepackage{fancyhdr}
\usepackage{multirow}
\usepackage{makecell}

\geometry{
  top=2.5cm,
  bottom=2.5cm,
  left=3cm,
  right=3cm,
}

%% ------------------------------------------------------------------
%% Colors and hyperref
%% ------------------------------------------------------------------
\definecolor{linkblue}{RGB}{37,99,235}
\definecolor{codeblue}{RGB}{99,102,241}
\definecolor{codebg}{RGB}{248,250,252}
\definecolor{best}{RGB}{22,163,74}
\definecolor{secbest}{RGB}{234,88,12}

\hypersetup{
  colorlinks=true,
  linkcolor=linkblue,
  urlcolor=linkblue,
  citecolor=linkblue,
  pdftitle={Week 14: Final Integrated Delivery},
  pdfauthor={Grupo 06 --- Big Data},
}

%% ------------------------------------------------------------------
%% Code listings
%% ------------------------------------------------------------------
\lstset{
  basicstyle=\ttfamily\footnotesize,
  backgroundcolor=\color{codebg},
  frame=single,
  rulecolor=\color{gray!30},
  breaklines=true,
  columns=flexible,
  keywordstyle=\color{codeblue}\bfseries,
  commentstyle=\color{gray},
  stringstyle=\color{teal},
}

%% ------------------------------------------------------------------
%% Header / footer
%% ------------------------------------------------------------------
\pagestyle{fancy}
\fancyhf{}
\rhead{Grupo 06 --- Big Data}
\lhead{Week 14: Final Integrated Delivery}
\rfoot{\thepage}

%% ------------------------------------------------------------------
%% Title
%% ------------------------------------------------------------------
\title{
  \vspace{-1cm}
  \Large\textbf{Week 14 Milestone Report}\\[0.4em]
  \large Final Integrated Delivery\\[0.3em]
  \normalsize MovieLens Discovery System --- Semester Project
}
\author{
  Grupo 06 --- Big Data\\
  \texttt{BD-Grupo-06/movielens-discovery}
}
\date{\today}

%% ==================================================================
\begin{document}
%% ==================================================================

\maketitle
\thispagestyle{fancy}
\tableofcontents
\newpage

%% ------------------------------------------------------------------
\section{Executive Summary}
%% ------------------------------------------------------------------

This report closes the semester arc for the MovieLens Discovery project. It does not introduce a
new technical layer; it integrates the five prior milestones (Week~3 ingestion, Week~5
representation, Week~7 clustering, Week~10 recommendation, Week~12 graph analytics) into one
defensible system, states how to rebuild it end to end, and is explicit about what should not be
overclaimed.

\begin{table}[H]
\centering
\caption{Headline result per layer}
\begin{tabular}{p{2.6cm}p{7.6cm}p{2.4cm}}
\toprule
\textbf{Layer} & \textbf{Headline result} & \textbf{Evidence} \\
\midrule
Data (Wk.~3) & 25,000,095 ratings, 62,423 movies, 1,093,351 tags cleaned; 99.74\% interaction sparsity & \texttt{reports/week03/} \\
Representation (Wk.~5$\to$7) & Autoencoder (13-dim) beats PCA (30-dim) $\sim$15$\times$ on validation MSE (0.0347 vs.\ 0.5083) & \texttt{reports/week05/}, \texttt{reports/week07/} \\
Clustering (Wk.~7) & K-means, $k=4$ on the AE-13 embedding, silhouette 0.178, 4 interpretable segments & \texttt{reports/week07/} \\
Recommendation (Wk.~10) & \texttt{content\_cosine} NDCG@10 = 0.9925 (genre protocol); \texttt{svd\_collaborative} Hit Rate@10 = 7.9\% (LOO protocol) --- protocols \emph{reverse} the ranking & \texttt{reports/week10/} \\
Graph (Wk.~12) & 13,176-node / 155,331-edge similarity graph; PageRank vs.\ popularity Spearman $\rho=0.313$ (weak) & \texttt{reports/week12/} \\
Demo & Interactive web app (\texttt{web/}): recommender search UI + 3D graph viewer, served from precomputed artifacts & \texttt{web/} \\
\bottomrule
\end{tabular}
\end{table}

\textbf{What this milestone adds beyond the weekly reports}: a single reproducibility path from raw
MovieLens files to every artifact in this repository (\texttt{RUNBOOK.md}), an
operationalization/monitoring plan for what would need to change if this ran as a real service, and
an honest, project-wide limitations section that supersedes the per-week ones.

\textbf{What should not be overclaimed}: this is an \emph{offline}, catalog-level system validated
with proxy metrics (genre overlap, leave-one-out history). It has never served a live user, has no
online A/B evidence, and the two Week~10 evaluation protocols disagree on which model is ``best''
--- the correct reading is that they measure different things, not that one protocol is wrong.

%% ------------------------------------------------------------------
\section{Problem Statement}
%% ------------------------------------------------------------------

Users facing the MovieLens catalog (62,423 movies) need a way to discover relevant titles without
relying on browsing an undifferentiated list. The product question, unchanged since Week~3, is:

\begin{quote}
  \textit{Given a movie a user already knows, which other movies should be surfaced next, why are
  they related, and what does the catalog's structure look like at a level above individual
  recommendations?}
\end{quote}

This decomposes into the four questions the semester layers were built to answer:
\begin{enumerate}[nosep]
  \item Which movies are similar in feature/behavior space? $\to$ Week~5 representation.
  \item Which movies belong to the same latent segment? $\to$ Week~7 clustering.
  \item Which movies should be recommended next, given a seed movie? $\to$ Week~10 recommendation.
  \item Which movies are structurally central in the similarity graph, and how does that differ
    from popularity? $\to$ Week~12 graph analytics.
\end{enumerate}

The system is a \textbf{catalog discovery and ranking tool}, not a personalized user-facing
production recommender: no login, no serving infrastructure, and no live feedback loop exist. That
framing matters for every claim made below.

%% ------------------------------------------------------------------
\section{Domain Context}
%% ------------------------------------------------------------------

Entertainment recommendation is a mature domain with well-known failure modes: popularity bias
(mainstream items dominate), cold start (new items/users have no signal), and filter bubbles
(over-narrow similarity). The project deliberately builds three independent views of the same
catalog --- content similarity (Week~5/7 embeddings), behavioral co-occurrence (Week~10 SVD), and
structural centrality (Week~12 graph) --- so that these failure modes can be \emph{compared against
each other} rather than accepted from a single model. The Week~10 popularity-vs-SVD-vs-content
comparison and the Week~12 popularity-vs-PageRank comparison ($\rho=0.313$) are the direct evidence
that this multi-view design surfaces different signals rather than three copies of the same
ranking.

%% ------------------------------------------------------------------
\section{Dataset Sources and Access Conditions}
%% ------------------------------------------------------------------

\begin{itemize}[nosep]
  \item \textbf{Source}: MovieLens 25M, GroupLens Research, University of Minnesota --- public
    research/educational release (\url{https://grouplens.org/datasets/movielens/25m/}).
  \item \textbf{Access method}: direct HTTPS download, no authentication, no scraping;
    \texttt{scripts/download\_dataset.sh} or \texttt{scripts/build\_week03\_pipeline.py} fetch the
    official \texttt{ml-25m.zip}.
  \item \textbf{Citation}: F. Maxwell Harper and Joseph A. Konstan. 2015. \emph{The MovieLens
    Datasets: History and Context.} ACM TiiS 5, 4: 19:1--19:19.
  \item \textbf{Auxiliary source}: TMDb API (\texttt{scripts/fetch\_posters.py}), used only to fetch
    poster images for the $\approx$400 movies sampled into the Week~12 web-viewer subgraph --- a
    presentation detail, not a modeling input. Requires a user-supplied \texttt{TMDB\_API\_KEY}; no
    key is committed to the repository.
  \item \textbf{Full source inventory and license text}: \texttt{reports/week03/}, Sections 2 and 7
    (unchanged since Week~3).
\end{itemize}

%% ------------------------------------------------------------------
\section{Schema and Data Dictionary}
%% ------------------------------------------------------------------

The canonical schema was fixed in Week~3 and has not changed. Six raw tables (\texttt{ratings},
\texttt{movies}, \texttt{tags}, \texttt{links}, \texttt{genome-scores}, \texttt{genome-tags}) are
cleaned into four processed tables under \texttt{data/processed/week03\_v1/}:

\begin{table}[H]
\centering
\caption{Processed table grain, size, and key}
\begin{tabular}{lp{3.6cm}rl}
\toprule
\textbf{Table} & \textbf{Grain} & \textbf{Rows} & \textbf{Key} \\
\midrule
\texttt{ratings\_clean.parquet} & one (user, movie) rating & 25,000,095 & (userId, movieId) \\
\texttt{movies\_catalog.parquet} & one movie & 62,423 & movieId \\
\texttt{tags\_clean.parquet} & one (user, movie, tag) event & 1,093,351 & none (event log) \\
\texttt{movie\_genres.parquet} & one (movie, genre) pair & 107,245 & (movieId, genre) \\
\bottomrule
\end{tabular}
\end{table}

Every downstream artifact (Week~5 feature matrix, Week~7 embeddings, Week~10 recommendation
tables, Week~12 graph) is keyed on \texttt{movieId} and traces back to this layer. Full
column-level dictionary: \texttt{reports/week03/}, Section 5.

%% ------------------------------------------------------------------
\section{Preprocessing and Feature Engineering}
%% ------------------------------------------------------------------

Week~5 built a $62{,}423 \times 45$ movie-level feature matrix from the Week~3 tables: 6 numeric
aggregates (log1p-scaled rating/tag/genre counts, z-scored release year, rating std), 19 binary
genre indicators, and 20 binary top-tag indicators. Two redundancy checks were run rather than
assumed: \texttt{tag\_event\_count\_log} was dropped (Pearson $r=0.985$ with
\texttt{unique\_tag\_count\_log}), while \texttt{genre\_count\_log} vs.\
\texttt{unique\_tag\_count\_log} was confirmed independent ($r=0.292$) and both were kept. Full
rationale: \texttt{reports/week05/}, Section 2.

%% ------------------------------------------------------------------
\section{Dimensionality and Representation Analysis}
%% ------------------------------------------------------------------

Two representations were built and compared, not just one:

\begin{itemize}[nosep]
  \item \textbf{PCA} (Week~5) on the 45-feature matrix: PC1+PC2 explain only 12.90\% of variance
    (no dominant axis); 30 components retain 80.96\% of variance at MSE$=0.1904$. PC1 reads as
    popularity/richness, PC2 as genre tone (action/thriller/horror vs.\ comedy).
  \item \textbf{Autoencoder} (Week~7): a nonlinear bottleneck swept from 2--30 latent dimensions,
    compared directly against PCA at each size. A 13-dimensional bottleneck was chosen: validation
    MSE 0.0347 vs.\ PCA's 0.5083 at the same size --- roughly 15$\times$ better reconstruction.
\end{itemize}

\textbf{Why the autoencoder became the production representation}: it is the input to every
downstream layer from Week~7 onward (clustering, \texttt{content\_cosine} recommendations, and
indirectly the Week~12 graph inherits the Week~10 SVD space, a separate but comparably-motivated
behavioral embedding). PCA is retained in the report as the linear baseline the autoencoder had to
beat, not as a discarded dead end --- the $\sim$15$\times$ MSE gap is the evidence for the more
complex model, not a because-we-could justification. What neither method proves: that the learned
axes correspond to anything a human would recognize as ground truth. PC1/PC2 read as
popularity/genre by \emph{inspection} of top loadings, which is suggestive, not a validated causal
factor.

%% ------------------------------------------------------------------
\section{Clustering Analysis}
%% ------------------------------------------------------------------

K-means was run on the 13-dimensional autoencoder embedding, sweeping $k=2$--20 and tracking
inertia and silhouette jointly rather than picking the metric-optimal $k=2$. \textbf{$k=4$} was
chosen: silhouette$=0.178$ (not high, but the tightest genre-behavior clusters obtainable without
collapsing into two catch-all buckets):

\begin{table}[H]
\centering
\caption{Cluster profile summary (13-dim autoencoder embedding, $k=4$)}
\begin{tabular}{rrrp{3.6cm}p{3.6cm}}
\toprule
\textbf{Cluster} & \textbf{Size} & \textbf{Share} & \textbf{Genre signal} & \textbf{Rating behavior} \\
\midrule
0 & 15,273 & 24.5\% & Thriller/crime/horror/action & $+540$ ratings vs.\ mean, slightly below-average score \\
1 & 32,747 & 52.5\% & Drama/comedy/romance (catch-all) & $-235$ ratings, $-0.109$ rating vs.\ mean \\
2 & 8,837  & 14.2\% & Animation/children/adventure/fantasy & $+148$ ratings, $+0.154$ rating vs.\ mean \\
3 & $\sim$4,566 & 7.3\% & Documentary (99.6\% purity) & $-334$ ratings, $+0.521$ rating vs.\ mean (niche/prestige) \\
\bottomrule
\end{tabular}
\end{table}

\textbf{Honest limitation carried forward from Week~7}: no density-based method (DBSCAN) was run
--- the Week~7 report explicitly scoped it out in favor of a deeper autoencoder-vs-PCA comparison.
This is a real gap against the Week~7 rubric's requirement for a justified density method, not
something to paper over in this integrated report: the clustering evidence in this project rests on
K-means alone, and cluster~1 (52.5\% of the catalog) is a diagnosed catch-all, not a validated
segment.

%% ------------------------------------------------------------------
\section{Recommendation or Ranking System}
%% ------------------------------------------------------------------

Three item-to-item systems were built and evaluated under \textbf{two independent protocols} that
disagree on the winner --- this disagreement is the most important finding of Week~10, not a
contradiction to be smoothed over:

\begin{table}[H]
\centering
\caption{Recommendation systems under two evaluation protocols}
\begin{tabular}{lrr}
\toprule
\textbf{System} & \textbf{Genre NDCG@10 (n=4,994)} & \textbf{LOO Hit Rate@10 (n=10,000)} \\
\midrule
\texttt{popularity\_global} (baseline) & 0.8390 & 4.65\% \\
\texttt{content\_cosine} (AE-13 cosine) & \textbf{0.9925} & 3.68\% \\
\texttt{svd\_collaborative} (SVD, $k=50$) & 0.9496 & \textbf{7.95\%} \\
\bottomrule
\end{tabular}
\end{table}

\begin{itemize}[nosep]
  \item Under the \textbf{genre-relevance protocol} (candidate pool = each system's own top-20,
    relevance = shares $\geq1$ genre with the query), \texttt{content\_cosine} wins because it
    directly optimizes for the same embedding space that genres are visibly encoded in.
  \item Under the \textbf{Leave-One-Out protocol} (real held-out user history, last-rated item as
    test), \texttt{svd\_collaborative} wins by a wide margin, nearly 2$\times$ the baseline and more
    than 2$\times$ \texttt{content\_cosine} --- evidence that behavioral co-occurrence predicts what
    a specific user actually watches next better than genre-coherent content similarity does.
  \item A \textbf{Recall@K measurement defect} was found and documented in the genre protocol
    (denominator was the system's own candidate pool, not the full catalog, inflating
    \texttt{content\_cosine} Recall@20 to 0.9994); the LOO protocol's Hit Rate@K is the corrected
    proxy. This is disclosed, not hidden, because a defense question about ``why is Recall so high''
    needs a real answer.
\end{itemize}

\textbf{Framing}: this is a ranking/discovery task (item$\to$item), not a personalized
user-to-item production recommender --- there is no serving path that consumes a live user
session.

%% ------------------------------------------------------------------
\section{Graph Analytics}
%% ------------------------------------------------------------------

The Week~12 graph reuses the Week~10 SVD item-factor space rather than introducing an unrelated
similarity metric: nodes are the 13,176 movies with $\geq50$ ratings (i.e., movies with an SVD item
factor); edges are undirected, weighted by cosine similarity, kept only among each node's top-$k=15$
neighbors above a minimum similarity of 0.5 (chosen via a joint sweep over $k \in \{5,10,15,20\}$
and threshold $\in \{0.3,0.5,0.7\}$, not hand-picked). The resulting graph has 155,331 edges,
0.179\% sparsity, one giant component covering 99.72\% of nodes.

\textbf{Key comparison finding}: PageRank correlates weakly with popularity (Spearman
$\rho=0.313$) and moderately with the Week~10 model's recommendation in-degree ($\rho=0.578$).
High-PageRank movies (e.g.\ \emph{L'Atalante}, \emph{Viridiana}) sit outside the popularity
top-1,000 entirely --- graph centrality answers ``how structurally embedded is this movie in the
similarity space,'' which is demonstrably not the same question as ``how many people rated it.''

\textbf{What the graph does not mean}: it is not a content/plot/cast similarity graph (built purely
from rating co-occurrence via SVD), it is not directed (so it cannot encode ``recommend $A$ because
of $B$'' --- that asymmetric signal lives separately in \texttt{model\_rec\_indegree}), and it
inherits whatever population bias exists in MovieLens raters.

%% ------------------------------------------------------------------
\section{Evaluation Protocol (Cross-Layer Summary)}
%% ------------------------------------------------------------------

\begin{table}[H]
\centering
\caption{Validation approach and acknowledged limitation, per layer}
\begin{tabular}{lp{5.3cm}p{5.3cm}}
\toprule
\textbf{Layer} & \textbf{Validation approach} & \textbf{Limitation acknowledged} \\
\midrule
Representation & AE vs.\ PCA validation-set MSE at matched latent size & Linear PCA baseline only; no comparison against sparse/topic-model text representations \\
Clustering & Silhouette + inertia sweep over $k=2$--20 & No density-based (DBSCAN) counter-check; silhouette$=0.178$ is modest in absolute terms \\
Recommendation & Genre-relevance protocol (proxy) \emph{and} LOO protocol (real user history) & The two protocols reverse the ranking; genre protocol had a documented Recall@K defect, now corrected via LOO Hit Rate@K \\
Graph & Sensitivity sweep ($k \times$ threshold) + two independent ranking baselines & Betweenness centrality is an approximate (500-source sample), not exact, at this graph size \\
\bottomrule
\end{tabular}
\end{table}

The throughline across all four layers: every major result is backed by a sweep or a second
comparison baseline, not a single hand-picked run. Where a sweep or baseline was \emph{not} done
(clustering's missing DBSCAN pass), it is named explicitly here and in Section~14 rather than
implied away.

%% ------------------------------------------------------------------
\section{Pipeline and Reproducibility}
%% ------------------------------------------------------------------

The full rebuild path --- from a fresh clone to every artifact referenced in this report --- is
documented in \textbf{\texttt{RUNBOOK.md}} at the repository root. Summary of what is and is not
script-driven today:

\begin{table}[H]
\centering
\caption{Reproducibility status per stage}
\begin{tabular}{p{3.6cm}p{5.4cm}p{4cm}}
\toprule
\textbf{Stage} & \textbf{Reproducible via} & \textbf{Status} \\
\midrule
Week~3 ingestion + cleaning & \texttt{scripts/build\_week03\_pipeline.py} & One-command, idempotent \\
Week~5 features + PCA & \texttt{scripts/build\_week05\_pipeline.py} & One-command, idempotent \\
Week~7 autoencoder + K-means & \texttt{notebooks/week07/*.ipynb} & Notebook-driven, no wrapper script \\
Week~10 recommenders + evaluation & \texttt{notebooks/week10/*.ipynb} & Notebook-driven, no wrapper script \\
Week~12 graph + subgraph export & \texttt{notebooks/week12/week12\_graph\_analytics.ipynb} & Notebook-driven, no wrapper script \\
Recommender catalog for demo & \texttt{scripts/build\_recommender\_catalog.py} & One-command \\
Poster URLs for graph viewer & \texttt{scripts/fetch\_posters.py} & One-command, optional \\
Web demo & \texttt{web/} (Astro + React) & \texttt{bun install \&\& bun run dev} \\
\bottomrule
\end{tabular}
\end{table}

\textbf{Honest gap}: Weeks 7, 10, and 12 are reproducible only by running notebooks top-to-bottom,
not through a single CLI command like Weeks 3 and 5 have. This is a real limitation against the
``hidden notebook state'' concern the rubric flags --- the notebooks were checked to run linearly
without relying on out-of-order execution, but they are not wrapped in \texttt{argparse}-driven
scripts the way \texttt{build\_week03\_pipeline.py} and \texttt{build\_week05\_pipeline.py} are.
\texttt{RUNBOOK.md} documents the exact run order so this gap does not block reproduction, but
closing it (extracting the notebooks into scripts) is listed as future work in Section~15.

%% ------------------------------------------------------------------
\section{Final Demo Artifact}
%% ------------------------------------------------------------------

The demo is the \texttt{web/} application (Astro~5 + React~19 + Tailwind), built directly on top of
already-produced, already-defended artifacts --- it computes nothing new:

\begin{itemize}[nosep]
  \item \textbf{Recommender view} (\texttt{/recommender}, \texttt{RecommenderView.tsx}): search the
    62,423-movie catalog by title (article-aware, e.g.\ ``Matrix'' matches ``Matrix, The (1999)''),
    pick a query movie, and see its top-10 \texttt{content\_cosine} and \texttt{svd\_collaborative}
    recommendations side by side, with posters, genres, and IMDb/TMDb links. Data comes from
    \texttt{artifacts/week10/week10\_recommender\_catalog.json}, built by
    \texttt{scripts/build\_recommender\_catalog.py} directly from the Week~10 evaluation artifacts.
  \item \textbf{Graph viewer} (\texttt{/}, \texttt{GraphViewer.tsx}): an interactive 3D
    force-directed rendering of the Week~12 sampled 400-node / 3,334-edge subgraph
    (\texttt{web/public/movie\_graph.json}), colored by Week~7 cluster, sized by PageRank, with
    hover/selection detail and poster art.
\end{itemize}

Both views are static-data consumers: the demo reads pre-materialized JSON, it does not call a
model at request time. That is a deliberate scope boundary (see Section~14) --- it demonstrates the
results of the pipeline, not a live inference service.

Run locally: \texttt{cd web \&\& bun install \&\& bun run dev}, then open
\url{http://localhost:4321}. See \texttt{RUNBOOK.md} Section~6 for the full command sequence,
including how to regenerate \texttt{movie\_graph.json} and \texttt{recommender\_catalog.json} if
the underlying week10/week12 artifacts change. Both routes were smoke-tested for this milestone: the
dev server serves \texttt{/} and \texttt{/recommender} with HTTP~200, and both the graph
(\texttt{movie\_graph.json}, $\sim$500~KB) and recommender catalog
(\texttt{recommender\_catalog.json}, $\sim$17~MB) load successfully in the browser.

%% ------------------------------------------------------------------
\section{Monitoring and Operationalization Plan}
%% ------------------------------------------------------------------

The system has never run as a service; this section states what \emph{would} need to exist if it
did, rather than describing monitoring that is actually deployed.

\textbf{If this were served in production:}

\begin{itemize}[nosep]
  \item \textbf{Serving assumption}: recommendation and graph artifacts would be precomputed in a
    batch job (matching the current notebook cadence) and served from a read-only store (the same
    JSON shape already consumed by \texttt{web/}, moved behind an API); no online training or
    per-request model inference is assumed.
  \item \textbf{Data drift}: MovieLens is a static academic snapshot, so there is no live rating
    stream to monitor in this project. If new ratings arrived, the metric to watch is the SVD
    explained-variance curve (currently 17.79\% at $k=50$, Week~10 Section~6.2) --- a sustained drop
    would indicate the latent space no longer explains new rating patterns well and a retrain is
    due. Sizing note: the Week~10 SVD training took $\sim$2.8s for $k=50$ on $162{,}540$ users
    $\times$ $13{,}176$ movies, so a full nightly retrain is cheap at this data scale.
  \item \textbf{Model drift}: the Week~10 genre-relevance vs.\ LOO disagreement (Section~8) is
    itself a warning that offline proxy metrics can diverge from what users actually do. In
    production, the metric to trust would be an online proxy of LOO (real held-out interactions),
    not the genre-overlap proxy, which this report shows is optimistic by construction for
    content-based methods.
  \item \textbf{Retraining cadence}: batch, on a schedule tied to how fast the catalog/ratings
    change --- for a static academic dataset this is moot, but for a live service the natural
    trigger is either a fixed interval (e.g.\ weekly SVD retrain) or a drift-metric threshold (SVD
    explained variance drop, recommendation-catalog coverage drop for the $\geq50$-rating floor
    used by both SVD and the Week~12 graph).
  \item \textbf{Logging}: none exists today. A production version would need to log query movie ID,
    returned candidate IDs, and (if available) subsequent user action, to make a real online
    evaluation possible --- the current LOO protocol is the offline stand-in for what that log
    would enable.
  \item \textbf{Failure modes to watch}:
  \begin{itemize}[nosep]
    \item \emph{Cold items}: any movie below the 50-rating floor has no SVD factor and no graph
      node (Week~10 documents 45,871 movies, 77.7\% of the catalog, excluded from the SVD matrix)
      --- these would silently fall back to content-only or popularity-only recommendations, and
      that fallback boundary is not currently visible to a user of the demo.
    \item \emph{Niche-genre blind spot}: \texttt{popularity\_global} systematically fails for
      Documentary, Western, Musical, Film-Noir (Week~10, Section~9.2) --- if popularity were ever
      used as a cold-start fallback in production, this failure mode would need an explicit
      override.
    \item \emph{Graph disconnection}: 29 isolated nodes and 33 total components at the chosen graph
      configuration (Week~12, Section~5) --- a monitoring check on isolated-node count would catch
      configuration drift if the graph were rebuilt with different data.
  \end{itemize}
\end{itemize}

%% ------------------------------------------------------------------
\section{Ethics and Limitations}
%% ------------------------------------------------------------------

\subsection{Ethics (unchanged since Week 3, reaffirmed here)}

\begin{itemize}[nosep]
  \item \textbf{Source and license}: GroupLens MovieLens~25M, public research/educational release,
    correctly attributed (Section~3).
  \item \textbf{Personal-data risk}: user IDs are anonymized integers; every artifact in this
    project (features, embeddings, clusters, recommendations, graph) operates at the \emph{movie}
    level. No individual rating history, user identifier, or re-identification attempt appears in
    any output, notebook, report figure, or the \texttt{web/} demo.
  \item \textbf{Mitigation carried through every layer}: aggregation at the movie level, no linkage
    to external personal data, raw data kept out of version control (\texttt{data/raw/},
    \texttt{env/} gitignored), and every transformation step documented for audit
    (\texttt{RUNBOOK.md}).
\end{itemize}

\subsection{Project-Wide Limitations (supersedes the per-week limitations sections)}

\begin{enumerate}[nosep]
  \item \textbf{Offline-only evaluation.} Every number in this report --- NDCG, Hit Rate,
    silhouette, PageRank correlation --- comes from an offline proxy (genre overlap, held-out
    history, static catalog structure). None of it has been checked against a live user. The LOO
    Hit Rate@10 of 7.95\% for the best system should be read as ``predicts held-out history 8\% of
    the time,'' not as ``will satisfy 8\% of real users'' --- those are different claims.
  \item \textbf{The two Week~10 protocols disagree, and that disagreement is the headline finding,
    not noise.} A team that reports only genre-relevance NDCG (favoring \texttt{content\_cosine})
    or only LOO Hit Rate (favoring \texttt{svd\_collaborative}) is not entitled to claim ``the best
    system'' without qualifying which question it answers.
  \item \textbf{Clustering rests on K-means alone.} No density-based method was run (Section~7);
    cluster~1 (52.5\% of the catalog) is a diagnosed catch-all bucket, not a clean segment, and the
    silhouette score (0.178) is modest in absolute terms.
  \item \textbf{Coverage gap at the low-rating tail.} 77.7\% of the catalog (45,871 movies) has no
    SVD factor and, by construction, no graph node --- the collaborative and graph layers only
    speak for the 22.3\% of movies with $\geq50$ ratings. Content-based recommendations
    (\texttt{content\_cosine}) are the only system that covers the full catalog.
  \item \textbf{Reproducibility is uneven across weeks.} Weeks 3 and 5 have single-command,
    argparse-driven pipelines; Weeks 7, 10, and 12 are reproducible only by running notebooks in a
    documented order (Section~11). This is a real, stated gap, not a hidden one.
  \item \textbf{The demo is a static-artifact viewer, not a live system.} It proves the pipeline's
    outputs are inspectable and interactive; it does not prove the system would perform acceptably
    under live traffic, concurrent users, or a changing catalog.
  \item \textbf{Single held-out split, no confidence intervals.} The LOO evaluation uses one random
    sample of 10,000 users at a fixed seed (42); no repeated resampling or variance estimate is
    reported, so the 7.95\% vs.\ 4.65\% vs.\ 3.68\% Hit Rate@10 gap is not accompanied by a
    significance test.
\end{enumerate}

%% ------------------------------------------------------------------
\section{Final Conclusions and Future Work}
%% ------------------------------------------------------------------

\subsection{What the semester built}

A movie discovery system that goes from 25M raw ratings to four independently-validated views of
the same 62,423-movie catalog --- content representation, behavioral segmentation, ranking, and
structural graph position --- plus a working interactive demo that renders the last two views
directly. Every major modeling choice (autoencoder latent size, K-means $k$, SVD rank, graph
$k$/threshold) is backed by a swept comparison rather than a single hand-picked run, and every layer
states a baseline it had to beat (PCA for the autoencoder, popularity for the recommenders,
popularity and model in-degree for the graph).

\subsection{Strongest result}

The Week~10 genre-relevance vs.\ LOO disagreement (Section~8): it is strong precisely because it is
inconvenient --- a team motivated to show one clean ``winning model'' would have reported only one
protocol. Reporting both, and explaining \emph{why} they disagree (proxy metric vs.\ real
behavioral signal), is the most defensible piece of analysis in the project.

\subsection{Weakest part, and what would be fixed first}

The clustering layer (Section~7 / Section~14.2): no density-based counter-check, a modest
silhouette score, and a 52.5\%-share catch-all cluster. With more time, the first fix would be a
DBSCAN or HDBSCAN pass on the same AE-13 embedding, explicitly to test whether cluster~1 is one
segment or an artifact of K-means forcing a spherical partition on non-spherical structure.

\subsection{What would change with more data, more time, or a different constraint}

\begin{itemize}[nosep]
  \item \textbf{More data} (a live rating stream instead of a static snapshot): the monitoring plan
    in Section~13 would become load-bearing rather than hypothetical --- SVD explained-variance
    drift and recommendation-catalog coverage would need real dashboards, not a written plan.
  \item \textbf{More time}: extract Weeks 7/10/12 notebooks into \texttt{argparse} pipeline scripts
    matching \texttt{build\_week03\_pipeline.py} / \texttt{build\_week05\_pipeline.py}, add a DBSCAN
    pass, and add confidence intervals to the LOO evaluation via repeated resampling.
  \item \textbf{A different constraint} (real user traffic instead of offline proxies): the
    genre-relevance protocol would be dropped in favor of a true online or A/B evaluation; the
    disagreement documented in Section~8 is exactly the reason an offline proxy alone should not
    gate a production launch decision.
\end{itemize}

The project is complete against its own product question: it can name, for any movie in the
62,423-movie catalog, what is similar to it (content), what is co-consumed with it (behavior), what
segment it belongs to (cluster), and how structurally central it is (graph) --- and it is honest
about which of those answers are validated against real user behavior (only the LOO recommendation
result) versus validated only against internal, proxy, or descriptive checks (everything else).

\end{document}
